\textbf{Implication for evaluation practice.} Symbolic and scripted
simulation is necessary but not sufficient for validating an LLM/VLM
robot planner. It validates reasoning strategy and task decomposition.
By construction it cannot exercise the failure modes of
Section~\ref{sec:failures}, because it never asks the model to compute
a quantity a physical sensor could contradict. We therefore
propose a \emph{geometry-grounding audit} as a standard step before
hardware deployment: identify every numeric quantity a prompt asks the
model to compute itself, and ask whether a sensor or a paired action
already determines it more reliably. Each quantity that fails the
audit is a candidate for a \harness{}-style correction.

\textbf{The cost of the architectural choice.} The related work argued
for a decoupled Sense-Plan-Act architecture over an end-to-end VLA on
data-efficiency and diagnosability grounds. The failure investigation
is the other half of that trade. \rev{The decoupled architecture asks
the model for physical quantities it was never trained to estimate,
and the six documented failure modes are the consequence; the harness
is the extra layer this architecture needs to be reliable in
practice.} We read this as a cost worth paying, not a case against the
architecture. \rev{Every resolved failure was fixable once diagnosed,}
precisely because the pipeline stayed decoupled enough to diagnose it.

\rev{\textbf{The unresolved collision case.} The sequential-placement
collision is a trajectory-level, multi-body contact problem, and the
static spacing mitigation of Section~\ref{sec:harness} does not
address its dynamic half. The most natural extension is an
obstacle-aware approach skill that treats already-placed objects as
static obstacles, optionally paired with a settle-time check after
each placement. This is the next addition we plan.}

\textbf{Generalizability.} Our evidence comes from one manipulator
(UR5cb), one gripper, one camera rig, and a single underlying model
(\code{kimi-k2.6}) used for both \rev{implementations}. We do not
claim that the six specific failures of Table~\ref{tab:taxonomy} recur
verbatim on different hardware. We claim that the structural pattern
is a property of the architecture, not of this robot: geometry
computed by a model instead of measured, and multiple real bodies
sharing space over time. Grasp width is the one quantity we still
cannot cross-check against depth in every case, since a single
top-down sample determines height unambiguously but not left-right
extent; the best-effort inset sampling of Section~\ref{sec:harness}
reduces this gap without \rev{eliminating} it.
